\section{Retrospective diagnostics for the holdout comparison}
\label{s:densitydiagnostics}

This appendix records corrections and sensitivity analyses performed after
the prospectively scheduled reads.  They test the frozen comparison on its
stated finite window; they do not alter the prospective rule or enlarge the
scope of the empirical claim.

\begin{remark}[Frozen readout and post hoc diagnostics]\label{r:sigma}
Three layers must be kept distinct.  First, the frozen inputs were the three
intensity functions, their numerical parameters, the baseline count, and---
before the first scheduled read---the reading times.  Second, the live tool
selected the closest displayed mean and printed Wald-type discrepancies
$(N-\mu)/\sqrt{N}$ for the bare and constant-floor candidates.  It did not
compute deviance residuals or exact Poisson tails; two manually noted
standardizations for the decaying candidate were later found to use the wrong
denominator.

Third, all deviance, dispersion, autocorrelation, negative-binomial, and block
calculations are retrospective.  The archival script used chunk midpoints as
the two endpoints of a cumulative integral.  Table~\ref{tab:milestones}
corrects that convention by integrating the full observed half-open chunk
intervals.  This changes no observation and no frozen operational intensity
function, only the integration limits.  The corrected means are therefore
post-campaign recalculations of the frozen functions.  Replacing the
operational $\widetilde\Lambda=p/3^{17}$ by the natural
$\Lambda=(p-1)/3^{17}$ changes every holdout total by less than
$9\cdot10^{-7}$ and changes no cell of Table~\ref{tab:milestones} at two
decimal places.
\end{remark}

\paragraph{Dispersion and dependence.}
On the $335$ post-freeze chunks, the frozen decaying model gives Pearson
statistic $X^2=324.2$ (quasi-Poisson scale $X^2/335=0.968$); an NB2 sensitivity
fit puts the additional-dispersion parameter at the Poisson boundary, and the
Poisson deviance is $210.46$.  The Pearson residuals have Ljung--Box
$Q_{12}=7.17$ (asymptotic conditional reference $p=0.85$), and aggregation into
blocks of $5,10,20$ chunks gives Pearson scales $0.81,0.96,1.05$.  These checks
reveal no gross overdispersion or serial correlation for that fixed model, but
they remain post hoc and have limited power in a sparse, non-stationary tail.

\paragraph{Block sensitivity of the forecast ordering.}
A moving-block resampling of the per-chunk log-score difference (circular
blocks of length $10$, $20\,000$ seeded replicates) gives a mean
decaying-minus-bare advantage of $0.01185$ per chunk, with a $95\%$ percentile interval
$[-0.00509,0.02957]$ and bootstrap fraction $0.092$ at or below zero.  For
decaying minus the frozen constant-floor model, the corresponding interval is
$[0.31298,0.49363]$ and no replicate is non-positive.  Thus the former
comparison is less decisive than the latter.  We report the forecast ordering,
not a universal fitted family or a discovery-level $p$-value.
